\begin{afterword}
\summaryhead

The essay uses Peter Wason's 2-4-6 task to describe a habit it calls positive bias. Subjects are told that 2-4-6 fits a hidden rule and may test other triplets. Most guess something like ``numbers increasing by two,'' test only triplets that fit their guess, hear ``yes,'' and announce the wrong rule. The real rule is any ascending sequence. Only about a fifth of subjects find it.

The author argues that this habit, testing cases your hypothesis predicts will fit rather than cases it predicts will fail, should be separated from confirmation bias in the sense of defending a belief you already hold. He links it to theories such as phlogiston that can explain any outcome. He holds that the habit works below the level of words, so that knowing about it does not cure it, and that people must train themselves to look for the cases that should not fit.

\responsehead

The essay teaches one lesson: look for evidence that would prove you wrong. It teaches it the way it tells the reader not to. It shows one task, designed by its inventor so that positive testing ``will almost certainly lead to erroneous conclusions,'' and draws from it a claim about ``human instinct.'' It does not look for the cases where its claim would fail. Those cases were in the literature for twenty years. Klayman and Ha showed in 1987 that testing expected positives is often the better strategy, and that it fails in 2-4-6 because the hidden rule was chosen to be broader than any natural guess. A post against positive bias argues from a single positive example.

The same paper also made the distinction the essay presents as its own: that the positive test strategy is not the same as defending one's belief. The essay offers it with ``it seems to me.'' The reader cannot tell that the one idea here that sounds new is not.

The essay's second claim, that the error is sub-verbal and cannot be fixed by knowing about it, is contradicted by the task's own research record. Asking subjects to find two complementary rules instead of one, a change of words, roughly triples the success rate. The essay reports the 20 per cent and not the 60.

The essay's examples drift. The student in its opening record did run a negative test, and learned nothing from it, because the hard part is choosing which negative test to run. The essay never addresses that problem. Phlogiston and ``emergence'' are about theories that explain everything, a different failure from choosing uninformative tests. The essay treats all three as one thing, and its advice, ``flinch toward the zero,'' gives no way to choose.

The essay offers its reader a humility that costs nothing. The author confesses that he too fell for the task, and invites the reader to confess the same. The last paragraphs then seal the thesis against its audience, so that nothing a reader does can count against it. The essay asks its reader to look into the dark and does not look there itself.

In short: a post against positive testing that tests its claim only on the one task built to confirm it, and presents a twenty-year-old distinction as its own.
\end{afterword}
